\section{Appendix: Scope and Qualifications}
\runninghead{Scope and Qualifications}

\subsection{Liturgical Identity, Occurrence, and Canon-Law Boundary}
This study examines the First Sunday of Advent (\latin{Dominica prima Adventus}, Sunday of the first class, \latin{Statio ad S.~Mariam maiorem}) in the 1962 typical edition of the \work{Missale Romanum} (\work{Proprium de tempore}, printed pp.~1--2; Preface of the Most Holy Trinity, n.~1082, printed p.~293), governed by the 1960 \work{Codex rubricarum} (\work{Acta Apostolicae Sedis}~52 [1960], pp.~597--729). Its dated occurrence is Sunday, 29~November~2026, under the universal General Roman Calendar, when the occurring ordinary commemoration of St.~Saturninus, Martyr, is omitted in Mass (RG~109, 111~a; RGMR~434) and the pre-1955 seasonal second and third orations remain abolished (1960 \work{Variationes}~II.13, IV.18). Particular diocesan, national, or religious calendars and annual Ordos (such as institutional papal prayers) govern local celebrations; present canonical authorization to celebrate Mass according to the 1962 Missal is a separate, mutable disciplinary matter outside this study's scope. Strict isolation from postconciliar calendars, readings, prayers, and cycles is maintained throughout.

\subsection{Text Verification, Translations, and Rights}
All ten proper texts, opening headings, and printed rubrics on pages~1--2 of the 1962 typical Missal were visually collated against the CMAA 1962 facsimile PDF (retained strictly as an unresolved-rights visual control, with no facsimile images reproduced). Because all ten Latin proper texts and the Preface of the Most Holy Trinity are preexisting pre-1931 liturgical wording corroborated word for word across tracked pre-1931 public-domain Missals (Pustet Ratisbon 1862, Tournai Descl\'ee 1889, Tours Mame 1922, and Cummiskey 1861), they are published under 17~U.S.C.~\S103(b) paired with those public-domain witnesses. Full collation notes and orthographic disclosures are recorded in \texttt{propers/verified.md}.

Historical English translations are included strictly for study and never as official 1962 liturgical English. Biblical texts follow the public-domain Douay--Rheims Bible (Challoner revision, 1899 American edition, in Vulgate psalm numbering), checked verse by verse against the Challoner Gutenberg 1581 witness (disclosing the single edition difference at Luke~21:29, ``in a similitude'' vs.\ ``a similitude,'' and avoiding Gutenberg's transcription error at Rom.~13:12). The three composed orations and the Preface of the Most Holy Trinity follow the public-domain historical translation of Fr.~F.~X. Lasance and Fr.~Francis Augustine Walsh, O.S.B., \work{The New Roman Missal in Latin and English} (Benziger Brothers, revised 1945, ordinary printed pp.~103, 105, 106, and 773--774, cleared as U.S.~public domain by failure to renew). All cited patristic, Doctoral, and medieval editions are U.S.~public domain, except the Wikisource Latin text of St.~Gregory the Great's \work{Homiliae in Evangelia} (licensed under CC~BY-SA~3.0, attributed in References alongside public-domain Migne PL~76, Bute 1908, and Oxford 1843).

\subsection{Author Standing, Liturgical Commentators, and Search Bounds}
Every whole-formulary interpretation rests on at least two distinct carrying authors registered in the repository's author-standing inventory (including at least one Father of the Church in each), with patristic exegesis of individual passages kept distinct from editorial synthesis across the complete Mass. St.~Bernard of Clairvaux (canonized 1174, Doctor of the Church 1830) is cited from verified tracked editions as a supporting Doctor; every interpretation already rests on at least two distinct carrying authors registered in the standing inventory. Amalarius of Metz, censured by a synod held at Kiersy (Quierzy) in 838 for his opinion concerning the signification of the parts of the divided Host at Mass, is cited solely for historical-structural witness to the Advent Epistle and \latin{Gloria} omission, never on the censured point. Medieval liturgical commentators such as Sicard of Cremona, Rupert of Deutz, and William Durandus report differing Advent~I Gospel assignments in some medieval uses, and are cited in this study solely for elements shared with the 1962 Mass.

Verified negative bounds and uninspected remote or restricted works are documented in \texttt{research/scope.md}: Cornelius a Lapide's Luke commentary omits verses~20--33 of chapter~21; St.~Thomas Aquinas's \work{Postilla super Psalmos} ends before Psalm~84; and tracked Aquinas Pauline OCR holds only 2~Corinthians through Colossians. Uninspected offline holdings include the restricted HTML payloads of the dated 2026 FSSP and ICRSP Ordos, the Latin text of RG~128 (violet vestment color is attested via registered FSSP Ordo metadata alongside verified RG~131 for Advent~III rose), the Murphy 1914 Douay--Rheims, Greek Psalmaries, Hilary of Poitiers, Cassiodorus on Psalms~24 and~84, St.~Augustine's \work{Epistula}~199 and \work{Expositio quarundam propositionum ex Epistola ad Romanos}, and St.~Thomas Aquinas's \work{Super Epistolam ad Romanos}. This document is a private study companion, not an official liturgical book, critical edition, or ecclesiastical approval.
